\scanpage{190}

$\underline{\overline{\mathrm{II}}}$\quad $\unsure{A|\mathcal A}(k^*)$. \deleted{\unsure{Le groupe}} \unsure{a)|1)|\textsuperscript{$\circ$})} Représentations irréductibles.

Se \unsure{fait|tout|font} via une représentation \unsure{rationelle|rationnelle} de $k^*$ dans $L^*$, où $L$ est un corps extension finie de $k$. $P(\lambda)=\sum_{-n}^{n}\lambda^i c_i$ $(c_i\in L)$

on a $P(\lambda\lambda')=P(\lambda)P(\lambda')$, i.e.

\[\sum_{\deleted{\unsure{i=-\infty|i}}}^{\unsure{\infty|n}}\lambda^i{\lambda'}^i c_i=\sum_{ij\deleted{\unsure{-\infty}}}^{\unsure{\infty|n}}\lambda^i{\lambda'}^j c_i c_j\unsure{=|,}\ ,\quad \text{et comme}\]

c'est une identité en $\lambda,\lambda'$ on doit avoir $c_i=c_i^2$, $c_i c_j=0$ si $i\ne j$, donc il ne peut y avoir qu'\deleted{\unsure{un des}} \uline{un} des $c_i$ qui soit $\ne 0$, soit $c_n$, et $\unsure{c=c^2|c_n=c_n^2}$ exige $c_n=1$. Donc $P(\lambda)=\lambda^n$ ; l'irréductibilité exige $L\unsure{:|=}k$. Ainsi

\margin{$k$ \uline{infini} (?)}

\noindent\begin{tabular}{@{}|p{.92\linewidth}@{}}\uline{Proposition} Les seules repr. \deleted{\unsure{idél.|irréd.}} irréductibles de $k^*$ ($k$ corps \uline{infini}) sont les repr. de dim 1, équivalentes aux repr. $\deleted{\unsure{\lambda\to\lambda^n}}\ \sigma_n(\lambda)t=\lambda^nt$.\end{tabular}

\uline{Corollaire} $\sigma_n\sigma_m=\sigma_{n+m}$, $\sigma_n^{\unsure{\prime|\vee}}=\sigma_{-n}$

Plus généralement, considérons une repr. lin. \inserted{\unsure{rationelle|rationnelle}} \uline{quelconque} de ${\unsure{\lambda|k}}^*$ dans un vectoriel $V$. On peut la mettre sous la forme

% archive sheet number: 189
